% Part 1 of Day 4: Number formats and the affine mapping.
%
% Topics of the syllabus: FP32, FP16, INT8; memory and energy per operation;
% scale, zero point, rounding error. Exercise: quantize five numbers by hand
% and calculate the error.
%
% Notation of this day: r is a real value, q is its integer, S is the scale,
% and Z is the zero point. The mapping is r = S (q - Z).

\theorypart{Number formats and the affine mapping}%
  {quantize five numbers by hand and calculate the error}

% ------------------------------------------------------------------- the cost
\begin{frame}{Each value of a model has a width}
  \small
  \renewcommand{\arraystretch}{1.3}
  \begin{tabular}{@{}lrrrr@{}}
    \toprule
    \textbf{Format} & \textbf{Bits} & \textbf{Bytes}
    & \textbf{Model of Day 3} & \textbf{MobileNetV2} \\
    \midrule
    FP32 & 32 & 4 & 25\,096 bytes & 13.4 MB \\
    FP16 & 16 & 2 & 12\,548 bytes & 6.7 MB \\
    INT8 & 8  & 1 & 6274 bytes    & 3.3 MB \\
    \bottomrule
  \end{tabular}
  \vskip 0.7em
  \begin{itemize}
    \item The table gives the bytes of the weights: the number of
          parameters times the bytes for each value.
    \item The model of Day 3 with the raw window has 6274 parameters.
          MobileNetV2 has 3\,504\,872 parameters.
    \item The same factor applies to the activations in RAM. The arena
          becomes smaller with the weights.
  \end{itemize}
  \vskip 0.2em
  \keypoint{Quantization changes the format of the values. It does not
    change the architecture: the layers and the number of parameters stay
    the same.}
\end{frame}

\begin{frame}{FP32: the format of the training}
  \begingroup\centering
  \begin{tikzpicture}[font=\footnotesize, x=0.31cm, y=0.8cm]
    \fill[cardinalred!40] (0,0) rectangle (1,1);
    \fill[darkcyan!35]    (1,0) rectangle (9,1);
    \fill[treegreen!35]   (9,0) rectangle (32,1);
    \draw[coolgray] (0,0) rectangle (32,1);
    \draw[coolgray] (1,0) -- (1,1) (9,0) -- (9,1);
    \node at (5,0.5) {exponent: 8 bits};
    \node at (20.5,0.5) {fraction: 23 bits};
    \node[anchor=north, font=\scriptsize] at (0.5,-0.05) {sign};
  \end{tikzpicture}\par\endgroup
  \vskip 0.4em
  \[
    r = (-1)^{\text{sign}} \times 1.\text{fraction} \times
        2^{\,\text{exponent} - 127}
  \]
  \begin{itemize}
    \item The exponent moves the binary point. So the format holds very
          small and very large numbers: from about $10^{-38}$ to $10^{38}$.
    \item The precision is about 7 decimal digits, for a small number and
          for a large number.
    \item A training needs this range: gradients can be very small.
  \end{itemize}
  \vskip 0.2em
  \keypoint{An inference needs much less. The weights and the activations
    of one layer are in a small range that you know after the training.}
\end{frame}

\begin{frame}{FP16: half the bits, the same idea}
  \begingroup\centering
  \begin{tikzpicture}[font=\footnotesize, x=0.31cm, y=0.8cm]
    \fill[cardinalred!40] (0,0) rectangle (1,1);
    \fill[darkcyan!35]    (1,0) rectangle (6,1);
    \fill[treegreen!35]   (6,0) rectangle (16,1);
    \draw[coolgray] (0,0) rectangle (16,1);
    \draw[coolgray] (1,0) -- (1,1) (6,0) -- (6,1);
    \node at (3.5,0.5) {5 bits};
    \node at (11,0.5) {fraction: 10 bits};
    \node[anchor=north, font=\scriptsize] at (0.5,-0.05) {sign};
    \node[anchor=north, font=\scriptsize] at (3.5,-0.05) {exponent};
  \end{tikzpicture}\par\endgroup
  \vskip 0.6em
  \begin{itemize}
    \item The largest number is 65\,504. The precision is about 3 decimal
          digits.
    \item The weights need half the bytes. The loss of accuracy is very
          small for almost all models.
    \item A microcontroller has no arithmetic for FP16. The runtime
          changes the weights to FP32 before it calculates.
    \item So FP16 saves flash. It saves no RAM and no time on a
          microcontroller.
  \end{itemize}
  \begin{block}{Example}
    The LiteRT file of the Day 3 lecture has 27\,020 bytes with FP32
    weights and 15\,248 bytes with FP16 weights.
  \end{block}
  \source{File sizes: results of the converter \texttt{onnx2tf} 1.29.24 on
    the work computer of the course.}
\end{frame}

\begin{frame}{INT8: 256 values and no exponent}
  \begingroup\centering
  \begin{tikzpicture}[font=\scriptsize, x=1cm, y=1cm]
    % floating point: dense near zero, sparse far from zero
    \node[anchor=west, font=\footnotesize] at (-5.2,1.35)
      {\textbf{Floating point:} the step grows with the number};
    \draw[coolgray] (-5,0.8) -- (5,0.8);
    \foreach \x in {0.08,0.16,0.24,0.32,0.4,0.56,0.72,0.88,1.04,1.36,1.68,
        2.0,2.32,2.96,3.6,4.24,4.88} {
      \draw[darkcyan] (\x,0.68) -- (\x,0.92);
      \draw[darkcyan] (-\x,0.68) -- (-\x,0.92);
    }
    \draw[darkcyan] (0,0.62) -- (0,0.98);
    % integer: uniform
    \node[anchor=west, font=\footnotesize] at (-5.2,0.0)
      {\textbf{Integer:} the same step for all numbers};
    \draw[coolgray] (-5,-0.55) -- (5,-0.55);
    \foreach \k in {-16,...,16} {
      \draw[cardinalred] (\k*0.305,-0.67) -- (\k*0.305,-0.43);
    }
    \node[below] at (-4.88,-0.7) {$-128$};
    \node[below] at (0,-0.7) {0};
    \node[below] at (4.88,-0.7) {127};
  \end{tikzpicture}\par\endgroup
  \vskip 0.4em
  \begin{itemize}
    \item An \code{int8} value is an integer from $-128$ to 127: 256
          values.
    \item The format has no exponent. The distance between two values is
          always 1.
    \item An integer alone cannot hold the weight 0.0173. You need a rule
          that connects the integers with real numbers.
  \end{itemize}
  \keypoint{The rule is the subject of this part: one scale and one zero
    point for each tensor.}
\end{frame}

% ---------------------------------------------------------------- why integer
\begin{frame}{Why integers: memory}
  \small
  \renewcommand{\arraystretch}{1.3}
  \begin{tabular}{@{}L{0.42\textwidth}rrr@{}}
    \toprule
    \textbf{Model of Day 3 with the raw window} & \textbf{FP32}
    & \textbf{INT8} & \textbf{Factor} \\
    \midrule
    Weights in the flash, in bytes & 25\,096 & 6274 & 4 \\
    Input tensor, 300 values, in bytes & 1200 & 300 & 4 \\
    Peak of the activations, in bytes & 1280 & 320 & 4 \\
    \bottomrule
  \end{tabular}
  \vskip 0.7em
  \begin{itemize}
    \item The flash holds the weights. The arena holds the activations. The
          two become four times smaller.
    \item The complete file and the complete arena decrease by a smaller
          factor: the description of the graph and the data of the
          interpreter do not change.
    \item On Day 1, MobileNetV2 in \code{float32} did not fit the flash of
          the XIAO. In \code{int8} the weights need 3.3 MB.
  \end{itemize}
  \vskip 0.2em
  \keypoint{For a microcontroller, the memory is the first reason for
    quantization. The speed is the second reason.}
\end{frame}

\begin{frame}{Why integers: energy and time}
  \begin{columns}[T]
    \begin{column}{0.44\textwidth}
      \footnotesize
      \renewcommand{\arraystretch}{1.3}
      \setlength{\tabcolsep}{4pt}
      \begin{tabular}{@{}llr@{}}
        \toprule
        \textbf{Operation} & \textbf{Width} & \textbf{Energy} \\
        \midrule
        Integer addition & 8 bits  & 1 \\
        Float addition   & 32 bits & 30 \\
        Read from memory & 64 bits & 40\,000 \\
        \bottomrule
      \end{tabular}
      \vskip 0.5em
      The unit is the energy of one integer addition with 8 bits.
    \end{column}
    \begin{column}{0.52\textwidth}
      \begin{itemize}
        \item An integer operation needs a smaller circuit than a float
              operation.
        \item To move data costs much more than to calculate. With
              \code{int8}, each read moves a quarter of the bytes.
        \item Many small microcontrollers have no circuit for float
              arithmetic. A float operation is then a software function.
      \end{itemize}
    \end{column}
  \end{columns}
  \vskip 0.5em
  \keypoint{The gain in time depends on the device and on the kernels. You
    measure it on the board in Part D of the lab.}
  \source{Energy ratios: \emph{Machine Learning Systems}, Volume I,
    chapter 10.}
\end{frame}

% ------------------------------------------------------------ affine mapping
\begin{frame}{The idea: a ruler with 256 marks}
  \begingroup\centering
  \begin{tikzpicture}[font=\scriptsize, x=0.92cm, y=1cm,
      >={Stealth[length=4pt]}]
    \draw[->, coolgray] (-0.3,1.5) -- (10.5,1.5) node[right] {$r$};
    \draw[->, coolgray] (-0.3,0) -- (10.5,0) node[right] {$q$};
    \foreach \x/\r/\q in {0/$-1.0$/$-128$, 2.5/0.0/$-64$, 10/3.0/127} {
      \draw[darkcyan, line width=0.9pt] (\x,1.38) -- (\x,1.62);
      \node[above] at (\x,1.62) {\r};
      \draw[cardinalred, line width=0.9pt] (\x,-0.12) -- (\x,0.12);
      \node[below] at (\x,-0.12) {\q};
      \draw[coolgray, dashed] (\x,1.38) -- (\x,0.12);
    }
    \foreach \k in {1,...,39} {
      \draw[cardinalred!60] (\k*0.25,-0.07) -- (\k*0.25,0.07);
    }
    \node[darkcyan] at (6.2,2.0) {real values: $r_{\min} = -1.0$, \;
      $r_{\max} = 3.0$};
    \node[cardinalred] at (6.2,-0.6) {integers: $-128$ to 127};
  \end{tikzpicture}\par\endgroup
  \vskip 0.4em
  \begin{itemize}
    \item You know the range of a tensor: its smallest value $r_{\min}$ and
          its largest value $r_{\max}$.
    \item You put the 256 integers on this range, with the same distance
          between them.
    \item Each real value goes to the nearest mark.
  \end{itemize}
  \vskip 0.2em
  \keypoint{Two numbers describe the ruler: the distance between two marks,
    and the integer that is at the real value 0.}
\end{frame}

\begin{frame}{The affine mapping}
  \begin{columns}[T]
    \begin{column}{0.50\textwidth}
      \begin{block}{From the integer to the real value}
        \[ r = S\,(q - Z) \]
      \end{block}
      \begin{block}{From the real value to the integer}
        \[ q = \text{clamp}\!\left(\text{round}\!\left(\frac{r}{S}\right)
               + Z,\ -128,\ 127\right) \]
      \end{block}
    \end{column}
    \begin{column}{0.46\textwidth}
      \small
      \begin{description}\itemsep 0.2em
        \item[$r$] The real value.
        \item[$q$] The integer that the device stores.
        \item[$S$] The scale: a positive real number. It is the real
          distance between two integers.
        \item[$Z$] The zero point: the integer that represents the real
          value 0.
        \item[clamp] Limits the result to the range of \code{int8}.
      \end{description}
    \end{column}
  \end{columns}
  \vskip 0.5em
  ``Affine'' means a multiplication and a shift. One tensor has one $S$ and
  one $Z$. The model file stores them with the tensor.
  \source{Mapping of Jacob et al., ``Quantization and Training of Neural
    Networks for Efficient Integer-Arithmetic-Only Inference'', 2018.}
\end{frame}

\begin{frame}{Derive the scale and the zero point}
  The range $[r_{\min}, r_{\max}]$ must go to the range $[-128, 127]$.
  \vskip 0.3em
  \begin{columns}[T]
    \begin{column}{0.47\textwidth}
      \begin{block}{Scale}
        255 steps cover the range:
        \[ S = \frac{r_{\max} - r_{\min}}{255} \]
      \end{block}
    \end{column}
    \begin{column}{0.49\textwidth}
      \begin{block}{Zero point}
        $r_{\min}$ must give $q = -128$:
        \[ Z = \text{round}\!\left(-128 - \frac{r_{\min}}{S}\right) \]
      \end{block}
    \end{column}
  \end{columns}
  \vskip 0.4em
  \textbf{Example:} $r_{\min} = -1.0$ and $r_{\max} = 3.0$.
  $S = 4.0 / 255 = 0.015686$, \; $r_{\min} / S = -63.75$, \;
  $Z = \text{round}(-128 + 63.75) = \text{round}(-64.25) = -64$.
  \begin{itemize}
    \item The zero point is an integer. So the real value 0 has an exact
          integer, with no rounding error.
    \item This is necessary: a padding with zeros and the output of a ReLU
          must be exactly 0.
  \end{itemize}
\end{frame}

\begin{frame}{Quantize one value, and go back}
  With $S = 0.015686$ and $Z = -64$:
  \vskip 0.4em
  \begingroup\centering
  \begin{tikzpicture}[font=\footnotesize, >={Stealth[length=5pt]},
      b/.style={draw=coolgray, rounded corners=3pt, fill=boxgray,
                minimum width=2.3cm, minimum height=0.9cm, align=center}]
    \node[b] (a) at (0,0)   {$r = 0.7$};
    \node[b, draw=cardinalred] (b) at (4.1,0) {$q = -19$};
    \node[b] (c) at (8.2,0) {$\hat{r} = 0.70588$};
    \draw[->] (a) -- node[above, font=\scriptsize] {quantize} (b);
    \draw[->] (b) -- node[above, font=\scriptsize] {dequantize} (c);
  \end{tikzpicture}\par\endgroup
  \vskip 0.4em
  \begin{columns}[T]
    \begin{column}{0.47\textwidth}
      \begin{block}{Quantize}
        $0.7 / 0.015686 = 44.63$\\
        round: 45\\
        $q = 45 + (-64) = -19$
      \end{block}
    \end{column}
    \begin{column}{0.47\textwidth}
      \begin{block}{Dequantize}
        $\hat{r} = 0.015686 \times (-19 + 64)$\\
        $\hat{r} = 0.015686 \times 45 = 0.70588$
      \end{block}
    \end{column}
  \end{columns}
  \vskip 0.5em
  \begin{itemize}
    \item The error is $\hat{r} - r = 0.00588$. The step is 0.015686.
    \item The device stores $-19$ in one byte. It does not store 0.7.
  \end{itemize}
  \keypoint{The mapping loses information. You cannot get 0.7 back from
    $-19$.}
\end{frame}

% --------------------------------------------------------------------- errors
\begin{frame}{The rounding error}
  \begin{columns}[c]
    \begin{column}{0.46\textwidth}
      \centering
      \begin{tikzpicture}[font=\scriptsize, x=0.62cm, y=0.62cm]
        \draw[->, coolgray] (0,0) -- (6.6,0) node[right] {$r$};
        \draw[->, coolgray] (0,0) -- (0,6.6) node[above] {$\hat{r}$};
        \draw[coolgray, dashed] (0,0) -- (6.2,6.2);
        \foreach \k in {0,...,5} {
          \draw[cardinalred, line width=1pt] (\k,\k+0.5) -- (\k+1,\k+0.5);
        }
        \foreach \k in {1,...,5} {
          \draw[cardinalred, line width=1pt] (\k,\k-0.5) -- (\k,\k+0.5);
        }
        \draw[<->] (3.5,3.5) -- node[right] {$S/2$} (3.5,4.0);
        \draw[<->] (2,1.9) -- node[below] {$S$} (3,1.9);
      \end{tikzpicture}
    \end{column}
    \begin{column}{0.52\textwidth}
      \begin{itemize}
        \item Inside the range, each value goes to the nearest mark.
        \item The error is between $-S/2$ and $+S/2$.
        \item The typical error is $S / \sqrt{12}$, about $0.29\,S$.
        \item One more bit makes $S$ half as large. The error becomes half
              as large also.
      \end{itemize}
    \end{column}
  \end{columns}
  \vskip 0.5em
  \begin{block}{Example}
    Range $[-1.0, 3.0]$ with 8 bits: $S = 0.0157$, largest error 0.0078.
    The same range with 4 bits: $S = 4 / 15 = 0.267$, largest error 0.133.
  \end{block}
  \keypoint{A small scale gives a small rounding error. The scale is small
    when the range is small.}
\end{frame}

\begin{frame}{The clipping error}
  \begingroup\centering
  \begin{tikzpicture}[font=\scriptsize, x=1.0cm, y=1.6cm]
    \draw[->, coolgray] (-0.2,0) -- (10.2,0) node[right] {$r$};
    \draw[darkcyan, fill=darkcyan!25]
      plot[domain=0:6, samples=60, smooth] (\x,{1.2*exp(-(\x-2.5)*(\x-2.5)/1.5)})
      -- (6,0) -- (0,0) -- cycle;
    \fill[cardinalred] (9.2,0.06) circle (2pt);
    \node[cardinalred, above] at (9.2,0.1) {one large value};
    \node[darkcyan] at (2.5,1.4) {almost all values};
    \draw[treegreen, line width=1pt] (0,-0.25) -- (5.2,-0.25);
    \node[treegreen, anchor=west] at (5.3,-0.25) {range A: small step,
      the large value is cut};
    \draw[cardinalred, line width=1pt] (0,-0.55) -- (9.2,-0.55);
    \node[cardinalred, anchor=north west] at (0,-0.6) {range B: no value is
      cut, the step is larger};
  \end{tikzpicture}\par\endgroup
  \vskip 0.3em
  \begin{itemize}
    \item A value outside the range goes to $-128$ or to 127. This is
          clipping. Its error has no limit.
    \item Range A cuts the large value, and all other values have a small
          rounding error.
    \item Range B keeps the large value, and all other values have a larger
          rounding error.
  \end{itemize}
  \vskip 0.2em
  \keypoint{The selection of the range is a trade-off between the two
    errors. Part 2 shows how a calibration set finds the range.}
\end{frame}

\begin{frame}{Symmetric and asymmetric ranges}
  \small
  \renewcommand{\arraystretch}{1.3}
  \begin{tabular}{@{}L{0.17\textwidth}L{0.36\textwidth}L{0.37\textwidth}@{}}
    \toprule
    & \textbf{Symmetric} & \textbf{Asymmetric} \\
    \midrule
    Real range   & $[-a, +a]$ & $[r_{\min}, r_{\max}]$ \\
    Zero point   & $Z = 0$ & Any integer from $-128$ to 127 \\
    Integers     & $-127$ to 127 & $-128$ to 127 \\
    Used for     & Weights & Activations \\
    Reason       & The weights are around 0. With $Z = 0$, the
                   multiplication is more simple.
                 & After a ReLU, all values are positive. A symmetric
                   range uses only half of the integers. \\
    \bottomrule
  \end{tabular}
  \vskip 0.7em
  \textbf{Example:} an activation after a ReLU with the range $[0, 6]$.
  \begin{itemize}
    \item Asymmetric: $S = 6 / 255 = 0.0235$ and $Z = -128$. All 256
          integers are in use.
    \item Symmetric: $S = 6 / 127 = 0.0472$. The integers below 0 are
          never in use, and the step is two times as large.
  \end{itemize}
  \source{Rules of the \texttt{int8} quantization specification of
    LiteRT.}
\end{frame}

% ---------------------------------------------------------- integer arithmetic
\begin{frame}{One layer with integers only}
  A fully connected layer calculates $y = \sum_i w_i\,x_i + b$. Each real
  value has its mapping: $w = S_w\,q_w$, \; $x = S_x (q_x - Z_x)$, \;
  $y = S_y (q_y - Z_y)$.
  \vskip 0.3em
  \begingroup\centering
  \begin{tikzpicture}[font=\footnotesize, >={Stealth[length=5pt]},
      b/.style={draw=coolgray, rounded corners=3pt, fill=boxgray,
                minimum width=2.25cm, minimum height=1.0cm, align=center}]
    \node[b] (a) at (0,0)    {Multiply and\\add in \code{int32}};
    \node[b] (b) at (2.75,0) {Add the bias\\\code{int32}};
    \node[b, draw=cardinalred] (c) at (5.5,0) {Multiply by\\$M$, and shift};
    \node[b] (d) at (8.25,0) {Add $Z_y$,\\clamp to \code{int8}};
    \foreach \a/\b in {a/b, b/c, c/d} { \draw[->] (\a) -- (\b); }
  \end{tikzpicture}\par\endgroup
  \vskip 0.3em
  \[
    q_y = Z_y + M \left( \sum_i q_{w,i}\,(q_{x,i} - Z_x) + q_b \right)
    \qquad M = \frac{S_w\,S_x}{S_y}
  \]
  \small
  \begin{itemize}
    \item The sum of many products of \code{int8} values needs more than 8
          bits. The kernel uses a sum with 32 bits.
    \item $M$ is a real number below 1. The converter stores it as an
          integer and a shift. The device then needs no float operation.
    \item The result goes back to 8 bits. The next layer gets \code{int8}
          values again.
  \end{itemize}
  \source{Jacob et al., 2018.}
\end{frame}

\begin{frame}{The bias has 32 bits}
  \begin{columns}[T]
    \begin{column}{0.50\textwidth}
      \begin{itemize}
        \item The bias adds to the sum of the products
              $q_w\,(q_x - Z_x)$.
        \item One unit of this sum is the real value $S_w \times S_x$.
        \item So the bias uses this scale, with the zero point 0:
              \[ S_b = S_w \times S_x \qquad q_b =
                 \text{round}\!\left(\frac{b}{S_b}\right) \]
        \item $S_b$ is very small. The integer $q_b$ is then large: it
              needs 32 bits.
      \end{itemize}
    \end{column}
    \begin{column}{0.46\textwidth}
      \begin{block}{The model of the Day 3 exercise}
        The first layer of \code{hello\_world.tflite}, as Netron shows it:
        \vskip 0.3em
        {\small
        \renewcommand{\arraystretch}{1.2}
        \begin{tabular}{@{}lr@{}}
          Input scale   & 0.024574 \\
          Weight scale  & 0.004224 \\
          Bias scale    & 0.00010381 \\
        \end{tabular}}
        \vskip 0.3em
        $0.024574 \times 0.004224 = 0.00010381$
      \end{block}
    \end{column}
  \end{columns}
  \vskip 0.5em
  \keypoint{This answers the question of Day 3: the weights have 8 bits,
    and the biases have 32 bits. A model has few biases, so they cost
    little flash.}
  \source{Scale values: the file \texttt{hello\_world.tflite} of
    \texttt{Labs/hardware/HW-02/}.}
\end{frame}

\begin{frame}[fragile]{What changes in the sketch}
\begin{lstlisting}[style=cpp]
TfLiteTensor* input = interpreter->input(0);
TfLiteTensor* output = interpreter->output(0);

// The model has int8 input. Quantize the real value x.
const float S_in = input->params.scale;
const int Z_in = input->params.zero_point;
input->data.int8[0] = (int8_t)(roundf(x / S_in) + Z_in);

interpreter->Invoke();

// The model has int8 output. Dequantize the result.
const float S_out = output->params.scale;
const int Z_out = output->params.zero_point;
const float y = S_out * (output->data.int8[0] - Z_out);
\end{lstlisting}
  \small
  \begin{itemize}
    \item The tensor holds its scale and its zero point. The sketch reads
          them from the model file.
    \item The sketch does the first step and the last step of the mapping.
          All steps between them are integer steps of the kernels.
    \item For a classifier, the largest output is the class. You can
          compare the \code{int8} outputs directly, with no dequantization.
  \end{itemize}
  \source{The code follows the sketch \texttt{tflm\_hello} of
    \texttt{Labs/hardware/HW-02/}.}
\end{frame}

\begin{frame}{The three formats in comparison}
  \small
  \renewcommand{\arraystretch}{1.3}
  \setlength{\tabcolsep}{4pt}
  \begin{tabular}{@{}L{0.20\textwidth}L{0.22\textwidth}L{0.22\textwidth}L{0.24\textwidth}@{}}
    \toprule
    & \textbf{FP32} & \textbf{FP16} & \textbf{INT8} \\
    \midrule
    Bytes for a value  & 4 & 2 & 1 \\
    Values             & About 4 billion & 65\,536 & 256 \\
    Step               & Grows with the number & Grows with the number
                       & Constant: the scale \\
    Flash              & Baseline & Half & A quarter \\
    RAM on a microcontroller & Baseline & No change & A quarter \\
    Arithmetic         & Float & Float, after a change to FP32 & Integer \\
    Loss of accuracy   & None & Very small & Small. You must measure it. \\
    \bottomrule
  \end{tabular}
  \vskip 0.7em
  \keypoint{The format \code{int8} is the standard for a microcontroller: it
    decreases the flash, the RAM, and the time. FP16 decreases only the
    flash.}
\end{frame}

% -------------------------------------------------------------------- exercise
\exerciseframe{quantize five numbers by hand}{%
  \small
  A tensor has the range $r_{\min} = -1.00$ and $r_{\max} = 1.55$. The
  target format is \code{int8}: from $-128$ to 127.
  \begin{enumerate}
    \item Calculate the scale $S$ and the zero point $Z$.
    \item Quantize these five values. Then dequantize them, and calculate
          the error $\hat{r} - r$.
          \[ 0.000 \qquad 0.500 \qquad -0.333 \qquad 1.234 \qquad 2.000 \]
    \item Which value has the largest error? Why?
    \item What is the largest rounding error for a value inside the range?
  \end{enumerate}
  \vskip 0.3em
  \[
    S = \frac{r_{\max} - r_{\min}}{255} \qquad
    Z = \text{round}\!\left(-128 - \frac{r_{\min}}{S}\right)
  \]
  \[
    q = \text{clamp}\!\left(\text{round}\!\left(\frac{r}{S}\right)
        + Z\right) \qquad \hat{r} = S\,(q - Z)
  \]
}

\answerframe{quantize five numbers by hand}{%
  \small
  \textbf{1.} $S = 2.55 / 255 = 0.01$ and
  $Z = \text{round}(-128 + 100) = -28$.
  \vskip 0.3em
  \textbf{2.} The table:\par
  \renewcommand{\arraystretch}{1.2}
  \begin{tabular}{@{}rrrrr@{}}
    \toprule
    $r$ & $\text{round}(r / S)$ & $q$ & $\hat{r}$ & \textbf{Error} \\
    \midrule
    0.000    & 0     & $-28$ & 0.00    & 0.000 \\
    0.500    & 50    & 22    & 0.50    & 0.000 \\
    $-0.333$ & $-33$ & $-61$ & $-0.33$ & 0.003 \\
    1.234    & 123   & 95    & 1.23    & $-0.004$ \\
    2.000    & 200   & 127 (clamp) & 1.55 & $-0.450$ \\
    \bottomrule
  \end{tabular}
  \vskip 0.4em
  \textbf{3.} The value 2.000. It is outside the range. The sum
  $200 - 28 = 172$ is above 127, so the clamp gives 127. This is the
  clipping error.
  \vskip 0.3em
  \textbf{4.} $S / 2 = 0.005$. The errors 0.003 and $-0.004$ are below this
  limit. The values 0.000 and 0.500 are exactly on a mark.
}
